\chapter{Glosario}
\label{app:man-glosario}

\begin{description}
    \item[ABI] Application Binary Interface. Contrato binario entre
    \filepath{libgraphmatch} (C++) y el resto del sistema. Versión
    chequeada al cargar.
    \item[Canonical OCSF] Formato estándar (OCSF v1.4) al que todo
    mapping traduce sus eventos. Lingua franca del ingestor.
    \item[Drift] Desviación estadística observada entre un batch
    nuevo y el perfil esperado del mapping actual. Dispara
    re-review.
    \item[DSL] Domain-Specific Language. El lenguaje de hipótesis
    de GraphHunter (ver cap.~\ref{sec:man-dsl}).
    \item[Entidad] Nodo del grafo. Tipada (\code{User},
    \code{Host}, \code{IP}, etc.).
    \item[FieldConfig] Objeto de datos que mapea campos de un log
    crudo a roles semánticos (\code{TIME}, \code{ACTOR},
    \code{TARGET}, ...) y tipos de entidad.
    \item[Fast lane] Carril determinístico de ingesta productiva.
    Ejecuta VRL compilado sin involucrar LLM.
    \item[FFI] Foreign Function Interface. Puente entre Rust y el
    código C++ de \filepath{libgraphmatch}.
    \item[GNN] Graph Neural Network. Modelo ML que clasifica paths
    por probable clase de amenaza.
    \item[Hipótesis] Descripción en DSL de un patrón a buscar en
    el grafo.
    \item[Hunt] Una ejecución de hipótesis contra el grafo. Su
    resultado son paths scoreados.
    \item[IOC] Indicator of Compromise. Entidad (IP, dominio, hash)
    marcada como indicador exportable.
    \item[k-simplicity] Cota estructural sobre la explosión
    combinatoria del matcher. \code{\{k=3\}} en una hipótesis.
    \item[Learning lane] Carril que observa el fast lane y dispara
    re-review cuando detecta drift.
    \item[Mapping library] Persistencia de mappings aprobados. Cada
    mapping tiene \code{mapping_hash} estable.
    \item[Matcher] Motor de búsqueda de patrones sobre el grafo.
    Implementación SIMD con equivalencia DFS validada por
    property-testing.
    \item[MCP] Model Context Protocol. Protocolo estándar para
    exponer tools a asistentes de IA.
    \item[Model gate] Patrón arquitectónico: la funcionalidad ML
    es opcional; su fallo no degrada el camino base.
    \item[ONNX] Open Neural Network Exchange. Formato de modelo ML
    que el GNN usa.
    \item[Path] Secuencia de nodos conectados por aristas,
    devuelta por un hunt.
    \item[Provenance] Metadata que estampa cada evento canonical
    con el \code{mapping_hash} que lo produjo. Permite
    trazabilidad.
    \item[Relación] Arista del grafo. Tipada (\code{Auth},
    \code{Connect}, \code{Spawn}, etc.).
    \item[RAII] Resource Acquisition Is Initialization. Patrón C++
    usado en los wrappers de \filepath{libgraphmatch} para
    gestión automática de memoria.
    \item[Sesión] Workspace aislado con su grafo, catálogo,
    hunts, tags y notas.
    \item[Shearing layers] Modelo arquitectónico que agrupa código
    por velocidad de cambio.
    \item[SIMD] Single Instruction Multiple Data. Vectorización
    usada en \filepath{libgraphmatch} para acelerar el matcher.
    \item[Slow lane] Carril offline donde un LLM propone mappings
    y un humano los aprueba.
    \item[Typestate] Patrón de tipado que codifica estados
    (\code{Building}, \code{Finalized}) como parámetros de tipo.
    Usado en \code{GmGraph} para prevenir uso incorrecto en
    compile-time.
    \item[VRL] Vector Remap Language. El lenguaje al que se
    compilan los mappings para ejecución determinística en el fast
    lane.
\end{description}
